\documentclass[a4paper, serif, 10pt]{moderncv}

%% moderncv
% color and style
\moderncvcolor{black}
\moderncvstyle{banking}

%% page layout/margins
\usepackage[top=2.5cm, bottom=2.5cm, left=1.5cm, right=1.5cm]{geometry}

\nopagenumbers{}

%% Babel config for Norwegian language
% ref:
% - https://latex3.github.io/babel/guides/locale-norwegian.html
\usepackage[norsk]{babel}

%% fontspec
\usepackage{fontspec}

\IfFontExistsTF{Linux Libertine}{
  \setmainfont[Ligatures={TeX, Common}, Numbers=OldStyle]{Linux Libertine}
}{}
\IfFontExistsTF{Linux Libertine O}{
  \setmainfont[Ligatures={TeX, Common}, Numbers=OldStyle]{Linux Libertine O}
}{}

%% CTeX
% CJK support, used to show CJK characters in the resume
%
% - heading=false: disable ctex heading
% - punct=kaiming: use kaiming punctuations styles for CJK
% - scheme=plain: use plain scheme, do not override \normalsize font size
% - space=auto: space settings for CJK characters
%
% CJK fonts are resolved per locale language, e.g. Noto Serif/Sans CJK JP
% for Japanese resumes, with fallback to the next region if not installed
%
% ref:
% - http://ctan.mirrorcatalogs.com/language/chinese/ctex/ctex.pdf
\usepackage[UTF8, heading=false, punct=kaiming, scheme=plain, space=auto]{ctex}

\IfFontExistsTF{Noto Serif CJK SC}{
  \setCJKmainfont{Noto Serif CJK SC}
}{}
\IfFontExistsTF{Noto Sans CJK SC}{
  \setCJKsansfont{Noto Sans CJK SC}
}{}

%% line spacing
\usepackage{setspace}
\setstretch{1.125}

%% URL styling - use normal text instead of monospace
\urlstyle{same}

%% Patch \httplink and \httpslink to handle full URLs
% The original moderncv macros always prepend http:// or https:// to URLs.
% This causes issues when the URL already has a protocol (e.g., https://example.com)
% resulting in malformed URLs like https://https://example.com.
% This patch checks if the URL already contains "://" and uses it directly if so.
% Note: We use \str_set:Nx (x-expansion) to expand macros like \@homepage first.
\ExplSyntaxOn
\str_new:N \g_yamlresume_url_str

\RenewDocumentCommand{\httpslink}{O{}m}{
  \str_set:Nx \g_yamlresume_url_str {#2}
  \str_if_in:NnTF \g_yamlresume_url_str {://}
    { \url{#2} }
    { \url{https://#2} }
}

\RenewDocumentCommand{\httplink}{O{}m}{
  \str_set:Nx \g_yamlresume_url_str {#2}
  \str_if_in:NnTF \g_yamlresume_url_str {://}
    { \url{#2} }
    { \url{http://#2} }
}
\ExplSyntaxOff

\name{Ingrid Johansen}{}
\title{Registrert sykepleier med erfaring fra akuttmottak}
\phone[mobile]{+47 912 34 567}
\email{ingrid.johansen@example.no}
\homepage{https://www.ingridjohansen.no}

\address{Storgata 42, Oslo, Oslo, Norge, 0182}{}{}

\extrainfo{{\small \faLinkedin}\ \href{https://www.linkedin.com/in/ingrid-johansen-rn}{@ingrid-johansen-rn}}

\begin{document}

\maketitle

\section{Grunnleggende informasjon}

\cvline{}{\begin{itemize}
\item Registrert sykepleier med bachelor i sykepleie og åtte års erfaring fra akuttmottak og medisinsk sengepost
\item Trygg på systematisk pasientvurdering, triage, medikamenthåndtering og tverrfaglig samarbeid
\item Erfaren med journalsystemene DIPS og Epic, og med å dokumentere i henhold til gjeldende rutiner
\item Opptatt av pasientsikkerhet, faglig utvikling og gode pasientforløp
\end{itemize}}
\section{Utdanning}

\cventry{aug. 2014–juni 2017}
        {Bachelor, Sykepleie, Poeng: 3,8}
        {OsloMet – storbyuniversitetet}
        {\url{https://www.oslomet.no}}
        {}
        {\begin{itemize}
\item Utviklet klinisk vurderingsevne gjennom praksis i medisinsk og kirurgisk sengepost
\item Fikk grunnleggende opplæring i legemiddelhåndtering, prosedyrer og pasientobservasjon
\item Arbeidet i tverrfaglige grupper med pasientsikkerhet og etiske problemstillinger som tema
\item Skrev bacheloroppgave om tidlig identifisering av sepsis i akuttmottak
\end{itemize}
\textbf{Kurs}: Anatomi og fysiologi, Farmakologi og legemiddelhåndtering, Akuttmedisin og intensivsykepleie, Psykisk helse og rus, Klinisk praksis i medisinsk sengepost, Klinisk praksis i kirurgisk sengepost, Sykepleieteori og etikk, Helseinformatikk og dokumentasjon}
\section{Arbeidserfaring}

\cventry{sep. 2020–Nå}
        {Registrert sykepleier, akuttmottak}
        {Oslo universitetssykehus HF}
        {\url{https://www.oslouniversitetssykehus.no}}
        {}
        {\begin{itemize}
\item Utfører systematisk triage og initiell pasientvurdering etter ABCDE-prinsippet
\item Observerer og behandler pasienter med akutte medisinske og kirurgiske tilstander
\item Administrerer og dokumenterer legemidler i DIPS i henhold til gjeldende prosedyrer
\item Samarbeider tett med leger, spesialsykepleiere og andre faggrupper i akutte forløp
\item Bidrar til simuleringstrening og opplæring av nyansatte og studenter
\end{itemize}
\textbf{Nøkkelord}: Akuttmedisin, Triage, Pasientsikkerhet, Tverrfaglig samarbeid}

\cventry{aug. 2017–aug. 2020}
        {Sykepleier, medisinsk sengepost}
        {Diakonhjemmet sykehus}
        {\url{https://www.diakonhjemmet.no}}
        {}
        {\begin{itemize}
\item Hadde ansvar for pasientobservasjon, medikamenthåndtering og dokumentasjon på medisinsk sengepost
\item Utførte prosedyrer som blodprøvetaking, kateterisering og sårbehandling
\item Deltok i tverrfaglige visittrunder og oppfølging av individuelle pasientforløp
\item Veiledet studenter og nyansatte i kliniske prosedyrer og journalsystem
\item Bidro til forbedringsarbeid knyttet til fallforebygging og ernæringskartlegging
\end{itemize}
\textbf{Nøkkelord}: Medisinsk sykepleie, Pasientobservasjon, Medikamenthåndtering, Dokumentasjon}
\section{Språk}

\cvline{Norsk}{Morsmål eller tospråklig nivå \hfill \textbf{Nøkkelord}: Morsmål, Skriftlig, Muntlig}
\cvline{Engelsk}{Fullt profesjonelt nivå \hfill \textbf{Nøkkelord}: Faglig kommunikasjon, Dokumentasjon, Presentasjoner}
\cvline{Spansk}{Begrenset arbeidskunnskap \hfill \textbf{Nøkkelord}: Pasientkommunikasjon, Grunnleggende}
\section{Ferdigheter}

\cvline{Klinisk sykepleie}{Ekspert \hfill \textbf{Nøkkelord}: Pasientvurdering, Medikamenthåndtering, Sårbehandling, Kateterisering, Blodprøvetaking}
\cvline{Akuttmedisin}{Avansert \hfill \textbf{Nøkkelord}: Triage, HLR, ABCDE-vurdering, Akuttmottak, Trauma}
\cvline{Kommunikasjon og samarbeid}{Avansert \hfill \textbf{Nøkkelord}: Pasientkommunikasjon, Tverrfaglig samarbeid, Dokumentasjon, Veiledning, Konflikthåndtering}
\cvline{Teknologi og journalsystemer}{Viderekommen \hfill \textbf{Nøkkelord}: DIPS, Epic, Helsenorge.no, Medikamentmodul, Pasientadministrasjon}
\section{Utmerkelser}

\cventry{des. 2023}
        {Oslo universitetssykehus HF}
        {Årets sykepleier 2023}
        {}
        {}
        {\begin{itemize}
\item Tildelt for fremragende klinisk vurderingsevne og tydelig pasientfokus i akuttmottak
\item Anerkjent for å ha bidratt til tryggere pasientforløp og godt tverrfaglig samarbeid
\end{itemize}}
\section{Sertifikater}

\cventry{mars 2023}
        {Norsk Resuscitasjonsråd}
        {HLR for helsepersonell}
        {\url{https://www.nrr.org}}
        {}
        {}

\cventry{jan. 2022}
        {DIPS AS}
        {Sertifisering i DIPS Arena}
        {\url{https://www.dips.com}}
        {}
        {}
\section{Publikasjoner}

\cventry{mai 2023}
        {Tidlig identifisering av sepsis i akuttmottak}
        {Sykepleien Fag}
        {\url{https://sykepleien.no/fag/artikkel/sepsis-akuttmottak}}
        {}
        {\begin{itemize}
\item Fagartikkel om systematisk screening og tidlig behandling av sepsis i akuttmottak
\item Oppsummerer erfaringer fra eget forbedringsarbeid og relevante faglige anbefalinger
\end{itemize}}
\section{Prosjekter}

\cventry{jan. 2022–des. 2022}
        {Et kvalitetsforbedringsprosjekt ved medisinsk sengepost for å redusere forekomsten av liggesår.}
        {Forbedringsprosjekt: Reduksjon av liggesår}
        {\url{https://www.oslouniversitetssykehus.no/forbedringsprosjekt-liggesar}}
        {}
        {\begin{itemize}
\item Kartla risiko og igangsatte systematiske hudvurderinger ved innleggelse
\item Innførte sjekklister og opplæring i trykkavlastning for pleiepersonell
\item Resultatene viste en tydelig reduksjon i antall registrerte liggesår
\end{itemize}
\textbf{Nøkkelord}: Kvalitetsforbedring, Pasientsikkerhet, Trykkavlastning, Hudvurdering}
\section{Interesser}

\cvline{Friluftsliv}{Fjellturer, Langrenn, Kajakk}
\cvline{Faglig utvikling}{Simuleringstrening, Fagartikler, Kurs}
\section{Frivillig arbeid}

\cventry{jan. 2018–Nå}
        {Førstehjelpsinstruktør}
        {Norsk Røde Kors}
        {\url{https://www.rodekors.no}}
        {}
        {\begin{itemize}
\item Holder førstehjelpskurs for frivillige og lokalsamfunn
\item Underviser i HLR, skadestedsledelse og praktisk førstehjelp
\item Bidrar til å styrke beredskapen og tryggheten i nærmiljøet
\end{itemize}}
    
\end{document}